\documentclass[10pt,a4paper]{amsart}
\usepackage[margin=19mm]{geometry}
\usepackage{amsmath,amssymb,mathtools}
\usepackage[scr=rsfs,cal=euler]{mathalfa}
\usepackage{xfrac}
\usepackage[unicode,hidelinks,bookmarks=false]{hyperref}
\newcommand{\ie}{i.e.\ }
\newcommand{\cf}{cf.\ }
\newcommand{\resp}{respectively\ }
\newcommand{\Subsection}{\subsection}
\providecommand{\nobreakdash}{\nobreak\hbox{-}}
\setlength{\emergencystretch}{2em}
\multlinegap0pt
\usepackage{bm}
\usepackage{bbm}
\usepackage{dsfont}
\usepackage{xspace}
\usepackage{cancel}
\usepackage{mathtools}
\usepackage{amsthm}
\makeatletter
\@ifundefined{thm}{\newtheorem{thm}{Thm}}{}
\@ifundefined{lem}{\newtheorem{lem}[thm]{Lemma}}{}
\@ifundefined{prop}{\newtheorem{prop}[thm]{Proposition}}{}
\makeatother
\title[Simply interpolating and Carleson sequences for Hardy spaces]{Simply interpolating and Carleson sequences for Hardy spaces in the polydisc}
\author{Nikolaos Chalmoukis, Alberto Dayan}
\date{Source: arXiv:2412.09099v3 (CC BY 4.0)}
\begin{document}
\maketitle

\noindent\emph{Source and licence.} Simply interpolating and Carleson sequences for Hardy spaces in the polydisc. Version arXiv:2412.09099v3 (CC BY 4.0), distributed under the Creative Commons Attribution 4.0 International licence, as recorded in the arXiv licence metadata for this exact version and shown on the abstract page https://arxiv.org/abs/2412.09099v3. Full rights notice: licenses/arxiv-2412.09099-CC-BY-4.0.txt.\par\smallskip

\noindent\textit{Editorial scope.} Counted unit: the Prop whose source label is \texttt{prop:bound\_inv} in the article (source0.tex lines 348--353, source label \texttt{prop:bound\_inv}), delivered together with the article's own complete original proof of it (source0.tex lines 355--405, one \texttt{proof} environment of 51 lines). No mathematics is written, repaired or re-derived: statement and proof are the article's own text line for line. Required context retained uncounted: prop:psi\_estimate (lines 319--322). Every definition, display and label of those lines is retained unchanged. Omitted from this package and not redistributed: 1--318, 323--347, 354--354, 406--646. Nothing inside a retained excerpt is omitted or altered: every retained body is delivered exactly as the article's lines are written, so no omission receipt of any kind accompanies this package. Citations: the retained excerpts carry no citation command at all, so no bibliography entry of the article is transcribed and no reference is redistributed. Reference adaptation: none was needed and none was made; every reference inside the delivered unit resolves inside it, so no unresolved reference or citation is produced. Printed slips kept literal: no printed word, symbol, hypothesis or number of the counted statement or of its proof was altered. Layout and class: the article's own class is \texttt{amsart} with packages including amsmath, amssymb, amsthm, xcolor, orcidlink, comment, hyperref; the wrapper is the lane's pinned \texttt{amsart} editorial wrapper at 10pt on A4, which restates the theorem-like environments the retained text needs and adds the article's own macro definitions that the retained text needs (none), with extra wrapper packages (bm, bbm, dsfont, xspace, cancel, mathtools). The delivered numbering is the unit's own. Assets: no retained span contains a figure environment, a graphics-inclusion command, a PDF-page inclusion, a table or a TikZ picture; no image file is copied into this package, the package declares no asset dependency and no image file is redistributed with it. Licence: version arXiv:2412.09099v3 of this article is distributed under the Creative Commons Attribution 4.0 International licence, as recorded in the arXiv licence metadata for this exact version and shown on the abstract page https://arxiv.org/abs/2412.09099v3. A full rights notice is delivered at licenses/arxiv-2412.09099-CC-BY-4.0.txt. Characters: every retained excerpt is pure ASCII, so the delivered engine, XeLaTeX with its default Latin Modern text font, reports no missing character.
\begin{lem} \label{prop:psi_estimate}
	For all $t>0$ there exists a constant $C=C(d, t)>0$ such that 
	\[ | \langle \psi_{z,t}, \psi_{w,t} \rangle| \leq C| \langle g_z , g_w \rangle |^{1+t} \]
\end{lem}

\begin{prop}
\label{prop:bound_inv}
	For all $\gamma, \Delta>0$ there exists a constant $C=C(\gamma,\Delta)>0$ such that, if $1\le p \le2$ and $ \Lambda \subseteq \mathbb{D}^2$ is a finite collection of points which is  $\gamma$  weakly separated and $\Delta$ column bounded  and also satisfies 
	\[ \Vert S_\lambda \Vert_2 = \text{constant on} \,\,\, \Lambda, \]
	then there exists a bounded linear operator $R: \ell^p \to H^p(\mathbb{D}^2)	 $ which is a right inverse of $T_\Lambda^p: H^p(\mathbb{D}^2) \to \ell^p $ and $\Vert R\Vert_{\ell^p\to H^p(\mathbb{D}^2)}\leq C. $   
\end{prop}

\begin{proof}
 Let $t>0$ to be determined later and notice that 
 \begin{equation}\label{eq:power_column_bound}
		\sum_{n:n\neq m} | \langle g_n, g_m \rangle|^{1+t} \leq (1-\gamma^2)^{\frac{t-1}{2}}\sum_{n:n\neq m} | \langle g_n, g_m \rangle|^2 \leq \Delta (1-\gamma^2)^{\frac{t-1}{2}}.  
	\end{equation}
	Since no confusion arises, we will continue to denote by $\ell^p$ the space $\mathbb{C}^N$, where $N$ is the number of points of the sequence $\Lambda$, equipped with the $\ell^p$ norm. 
	Then we consider the operator $B_t:\ell^p \to H^p(\mathbb{D}^2)	$ defined as 
	\[ B_t(a) = \sum_{n=1}^{N} a_n \psi_{n,t}, \,\, a=(a_1,\dots,a_N)\in \ell^p,  \]
where $ \psi_{n,t}= \psi_{\lambda_n,t}.$ 

Let us start by estimating $\Vert B_t\Vert_{\ell^2 \to H^2(\mathbb{D}^2)} $. We apply first Lemma \ref{prop:psi_estimate} and subsequently Cauchy-Schwarz and the estimate \eqref{eq:power_column_bound} in order to get that, for some $C$ depending on $t$, we have 
\begin{align*} \Vert B_t(a)\Vert^2_{2}  & = \sum_{n,m=1}^N a_n \overline{a_m} \langle \psi_{n,t}, \psi_{m,t} \rangle \\
	& \leq C \sum_{n,m=1}^N |a_n a_m| |\langle g_n, g_m \rangle|^{1+t} \\
	& \leq C \sup_{m} \Big( \sum_{n=1}^{\infty} | \langle g_n, g_m \rangle|^{1+t} \Big)  \Vert a\Vert_{\ell^2}^2 \\ 
	& \leq C(1+\Delta (1-\gamma^2)^{\frac{t-1}{2}}) \Vert a\Vert_{\ell^2}^2.
 \end{align*}

 In order to estimate  $\| B_t\|_{\ell^1 \to H^1(\mathbb{D}^2)}$ set $s:=p/2(1+t)-1$ and notice that, $s>0$ for all $t>1$ and that


\begin{equation}
\begin{split}
\label{eq:pnorm_psit}\Vert \psi_{z,t} \Vert_p^p  &=\|\psi_{z, t}^\frac{p}{2}\|^2_2\\
&=\left\|\frac{S_z^{1+s}}{\|S_z\|_2^{\frac{p}{2}(1+2t)}}\right\|_2^2\\
&=\|\psi_{z, s}\|^2_2\cdot\|S_z\|^{-p(1+2t)+2+4s}\\
&\leq C\|S_z\|_2^{p-2}
\end{split}
\end{equation}

thanks to Lemma \ref{prop:psi_estimate}, $z=w$. In particular for $a\in \ell^1$
\[ \Vert B_t(a)\Vert_{\ell^1} \leq \Vert a\Vert_{\ell^1} \sup_{1\leq n \leq N} \Vert \psi_{n,t}\Vert_{1}  \leq C \|S_{\lambda_1}\|^{-1}_2 \Vert a\Vert_{\ell^1},  \]
since $\|S_{\lambda_n}\|_2$ is constant on $\Lambda$. Notice that by considering $H^p(\mathbb{D}^2)$ as a subspace of $L^p(\mathbb{T}^2,m)$ we can apply the Riesz-Thorin interpolation theorem to conclude that $\Vert B_t\Vert_{\ell^p \to H^p(\mathbb{D}^2)} \leq C\|S_{\lambda_1}\|_2^{1-\frac{2}{p}}$, where $C$ depends only on $\Delta, \gamma$ and $t.$   

Next we consider the operator $T_\Lambda^p B_t : \ell^p \to \ell^p$, which for $a\in \ell^p$ acts as follows 
\begin{align*}
	T_\Lambda^p B_t (a) & = T_\Lambda^p \Big ( \sum_{n=1}^N a_n \psi_{n,t} \Big) \\
			    & = \Big( \sum_{n=1}^N a_n \psi_{n,t}(\lambda_m) \Vert S_{\lambda_m}\Vert_2^{-\frac2p} \Big)_{m=1}^N  \\
			    & =  \Big( \sum_{n=1}^N a_n \frac{S_{\lambda_n}(\lambda_m)^{1+t}}{\Vert S_{\lambda_n}\Vert_2^{1+2t+\frac2p}}   \Big)_{m=1}^N \\
			    & =  \Big( \Vert S_{\lambda_1}\Vert_2^{1-\frac2p} \sum_{n=1}^N a_n \langle g_n, g_m \rangle^{1+t}  \Big)_{m=1}^N,
\end{align*}
where we have used the fact that the norm of the Szeg\"o kernel vectors is constant on the sequence.  

Therefore we conclude that 
\[
\Vert \Vert S_{\lambda_1}\Vert_2^{\frac2p-1} T_\Lambda^p B_t - Id \Vert_{\ell^p\to \ell^p } \leq \|(\langle g_n, g_m\rangle^{1+t})_{n, m=1}^N\|_{\ell^p\to\ell^p}\leq \Delta(1-\gamma^2)^{\frac{t-1}{2}} 
\]
via an applicarion of \eqref{eq:power_column_bound} and Riesz's interpolation theorem.
Choosing $t>0$ such that the latter expression equals $ \frac12 $ we conclude that the operator $T_\Lambda^pB_t $ is invertible on $\ell^p $. Moreover, the norm of its inverse is controlled by $\|S_{\lambda_1}\|^{\frac{2}{p}-1}_{2}$, up to a constant that is independent of $p$. Hence $R : = B_t(T_\Lambda^p B_t)^{-1} : \ell^p \to H^p(\mathbb{D}^2) $ is a right inverse of $T_\Lambda^p  $ and
\[ \Vert R \Vert_{\ell^p \to H^p(\mathbb{D}^2)} \leq \Vert B_t\Vert_{\ell^p \to H^p(\mathbb{D}^2)} \Vert (T_\Lambda B_t)^{-1}\Vert_{\ell^p\to \ell^p}\leq C, \]
where $C$ does not depend on $p$. This concludes the proof.
\end{proof}

\end{document}
